\section{模型评价与改进}

\subsection{模型优点}

按模型而不是按问题罗列优点；每项对应已完成的验证、数值或结构事实，不写“先进”“适用广泛”等空泛判断。

\subsection{模型局限}

对每项局限写清来源、可能造成的偏差和受影响的结论，例如样本跨度、关键变量缺失、参数漂移、简化约束或局部最优；不得把未来工作伪装成已完成结果。

\subsection{改进与推广}

区分已验证改进与未来建议。推广时说明哪些模型结构可复用，哪些参数、阈值、权重、约束和数据必须重新获取、校准和验证。
